\section{System Architecture}
\label{sec:architecture}

The architecture of \modora is designed to transform unstructured, multimodal documents into a structured knowledge hierarchy while providing an interactive and verifiable user experience. As illustrated in Figure~\ref{fig:arch}, the system is vertically organized into three layers: the \textbf{User Interaction Layer}, the \textbf{\modora Service Engine}, and the \textbf{Three-Stage Core Pipeline}.

\begin{figure*}[tp]
  \centering
  \includegraphics[width=0.9\linewidth]{Fig/arch(5).pdf}
  \caption{System Architecture of \modora.}
  \label{fig:arch}
\end{figure*}

\subsection{User Interaction Layer}
To bridge natural language queries with structured document trees, \modora provides an integrated environment featuring three specialized modules:

\runin{File Manager} Handles multimodal document ingestion and semantic tagging. \xbr{A persistent global library enables cross-session document reuse, eliminating redundant preprocessing.}%It maintains a persistent global library, enabling cross-session document reuse and eliminating redundant preprocessing.

\runin{Chat Engine} Supports complex QA with \textit{Interactive Grounding}. This feature allows users to trace AI-generated answers back to original PDF layouts via automated coordinate-based highlighting, ensuring the reasoning process is verifiable.

\runin{Tree Visualizer} Facilitates human-in-the-loop intervention by allowing users to manually refine the CCTree structure. It incorporates a \textit{Node Heatmap} (based on retrieval impact) to visually identify critical information hubs within the document hierarchy.

\subsection{\modora Service Engine}
Acting as the middleware bridging the user interaction layer and the core algorithmic pipeline, the Service Engine is optimized for multimodal RAG workloads. It orchestrates the asynchronous execution of computationally intensive VLM preprocessing tasks (Stage 1) and manages global document indexing to support cross-document question-answering across various sessions.

\subsection{Stage 1: Multimodal Preprocessing}
Before indexing, documents undergo a rigorous enrichment process to capture multimodal nuances:

\runin{Local-Alignment Aggregation} The system first performs optical character recognition (OCR) using \textit{PPStructureV3}~\cite{PPStructureV3} to identify raw elements such as titles, text blocks, tables, and images. Subsequently, a local-alignment strategy is applied using heuristic rules (e.g., text aggregation rule $A_1$, non-text aggregation rule $A_2$) to group fragmented elements into self-contained Aggregated Components. This ensures the components \xbr{are semantically (e.g., a table with its corresponding title) and visually (pixel-level coordinate) complete}.

\runin{VLM-based Information Enrichment} For non-textual components, \modora utilizes VLMs to extract integrated information, forming a semantic triplet consisting of Title, Metadata, and Data. For charts, the focus is on axes, labels, and trends; for tables, the schema and representative cell values are prioritized.

\subsection{Stage 2: Recursive Tree Construction}
The system organizes the enriched components into a logical hierarchy through a two-stage process:

\runin{Hierarchical Relationship Modeling} 1) \textit{Main Branch Construction}: Establishes ``text-to-text'' hierarchical relations based on title level prediction and ``text-to-other'' associations by attaching adjacent non-text components as children of relevant text nodes. 2) \textit{Supplementary Branch Isolation}: Page headers, footers, and sidebars are isolated into an independent branch under a virtual root to prevent semantic interference \xbr{with the main body} during retrieval.

\begin{sloppypar}
\runin{Bottom-up Metadata Summarization} To address information loss in long-document retrieval, summaries are recursively generated from leaf nodes upward. To balance information density, the number of summary keywords $n_i$ for each node $v_i$ is determined by an information attenuation formula, i.e., $n_i = \left\lceil \frac{D_i + d_i - 1}{D_i} \cdot \log_2(n_{i-1}^k) + 1 \right\rceil$, where $D_i$ denotes the maximum depth of the subtree rooted at $v_i$, $d_i$ represents the depth of node $v_i$, and $n_{i-1}$ indicates the total number of summaries generated by the children of $v_i$. The term $\frac{D_i + d_i - 1}{D_i}$ serves as the information attenuation coefficient, while $k$ is a hyperparameter controlling the rate of information growth. This ensures the root node captures a global overview while preserving critical details from the leaves.
\end{sloppypar}

\subsection{Stage 3: Two-Pathway Retrieval}
\modora uses a question-type-aware ``Parse-Retrieve-Verify'' workflow to balance recall and precision.

\runin{Query Parsing} The system first identifies whether a query contains spatial cues, then routes it to location-aware or semantic retrieval.

\runin{Location-Aware Path} For position-sensitive queries, \textit{LocationRetriever} traverses the CCTree with normalized bounding boxes and returns nodes overlapping the target regions, enabling direct grounding to layout coordinates.

\runin{Semantic Path} For content-centric queries, \textit{SemanticRetriever} recursively prunes irrelevant branches with LLM guidance, applies embedding fallback on pruned branches, and verifies candidate nodes using both text and cropped visual regions.

\runin{Verification-then-Fallback} Retrieved evidence is used for final reasoning. If consistency checks fail, the system falls back to whole-page reasoning. The retrieved nodes' impact scores are then updated for future tree optimization.
